\documentclass[conference]{IEEEtran}
\IEEEoverridecommandlockouts

% -----------------------------------------------------------------------------
% PACKAGES
% -----------------------------------------------------------------------------
\usepackage{cite}
\usepackage{amsmath,amssymb,amsfonts}
\usepackage{algorithmic}
\usepackage{algorithm}
\usepackage{graphicx}
\usepackage{textcomp}
\usepackage{xcolor}
\usepackage{hyperref}
\usepackage{booktabs}
\usepackage{multirow}
\usepackage{tikz}
\usetikzlibrary{arrows, positioning, shapes}

% -----------------------------------------------------------------------------
% ANONYMITY SETTINGS FOR PDF METADATA
% -----------------------------------------------------------------------------
\hypersetup{
    pdfauthor={Quhao Chen},
    pdftitle={Quhao_Chen},
    pdfsubject={HotShift},
    pdfkeywords={Machine Learning},
    colorlinks=true,
    linkcolor=black,
    citecolor=black,
    filecolor=black,
    urlcolor=blue
}

\begin{document}

% -----------------------------------------------------------------------------
% TITLE AND AUTHORS
% -----------------------------------------------------------------------------
\title{HotShift: Distilling Dense Reasoning from Proprietary Thought Summaries via Pivot-Token Reinforcement Learning}

\author{
\IEEEauthorblockN{Quhao Chen}
\IEEEauthorblockA{High School Student\\
Shanghai High School International Division\\
Shanghai, China\\
Email: willcxd@outlook.com}
}

\maketitle

% -----------------------------------------------------------------------------
% ABSTRACT AND KEYWORDS
% -----------------------------------------------------------------------------
\begin{abstract}
Small reasoning models are verbose, and verbosity is the tax that keeps them off small hardware. A stock Qwen3.5-4B in thinking mode emits $\sim$21K output tokens per long-horizon software-engineering task and still fails most of them; its forerunners at 4B resolve $\sim$11\% of SWE-bench Verified, and vanilla outcome-level GRPO moves that number by less than half a point. Meanwhile, proprietary teacher models worth distilling from increasingly obfuscate their chains of thought, exposing only a summary generated by an external model and a final answer. We present \textbf{HotShift}, a three-stage recipe that trains a 4B student to \textit{think in pivots}: (1) pivot-dense SFT, in which the student synthesizes its own compact token traces conditioned only on a proprietary teacher's summary ``bubbles'', discarding the low-entropy filler that dominates teacher rationales; (2) summary-hint step optimization, which localizes the single pivotal step where a failed trajectory diverges, queries the teacher for one abstract hint at that step only, and trains on outcome-verified failure$\to$success flips with credit concentrated at the pivot; and (3) dense-thinking GRPO, a group-relative objective with a joint correctness and length reward that teaches the stopping point SFT never learns. The entire pipeline --- 3.1K distilled traces, 1.8K verified pivot pairs, 300 GRPO steps --- runs on one NVIDIA T4 (16 GB) (Amazon SageMaker) with FP16 LoRA, at a total program cost of US \$23.80 (for API access; T4 compute is subsidized via free sessions). On a decisively selected 250-task portion of the contamination-free DeepSWE benchmark --- chosen for representativeness across repository, language, and difficulty strata, not for convenience --- results: pass@1 7.9\% $\to$ 19.4\% (+11.5 points, 2.5$\times$) while cutting mean reasoning tokens $-$78\% (21.4K $\to$ 4.7K), a 9.1$\times$ reduction in tokens per solved task. At a token budget matched to its own output length, the stock model collapses to 4.3\% while HotShift holds 19.4\%.
\end{abstract}

\begin{IEEEkeywords}
Reinforcement learning, Large language models, Software engineering agents
\end{IEEEkeywords}

% -----------------------------------------------------------------------------
% MAIN CONTENT
% -----------------------------------------------------------------------------
\section{Introduction}
\label{sec:introduction}
Three observations motivate HotShift. 

First, small reasoners overthink. Post-SFT small models ``fail to determine the optimal stopping point of the reasoning process, resulting in verbose and repetitive outputs'' \cite{tldr}. Qwen3.5-4B --- the strongest sub-5B reasoning model \cite{qwen35} --- burns 200M+ output tokens to complete a general intelligence index \cite{artificialanalysis}, and its maintainers recommend $\ge$128K context ``to preserve thinking capabilities'' \cite{qwen35}. On long-horizon software engineering tasks this verbosity is not a cost problem; it is a feasibility problem: 21K tokens per attempt at 300 tok/s (a T4's generation rate \cite{unsloth}) is a minute per rollout before the agent has done anything. Furthermore, in most local AI deployments, the model is frequently ran in quantized weights, which results in even more tokens drained during reasoning.

Second, most of those tokens do not matter. Only $\sim$20\% of CoT tokens are high-entropy ``forking'' tokens that steer the reasoning; RLVR restricted to those 20\% matches full training at 8B and beats it at 14B--32B, while training on the low-entropy 80\% degrades performance \cite{wang2025}. Reasoning failures localize to one or two pivotal steps --- errors in the first step alone cause a $\sim$40\% accuracy drop \cite{liao2025}, and critical-step optimization needs supervision at only $\sim$16\% of trajectory steps to beat full-trajectory SFT/DPO \cite{cso}.

Third, the good teachers no longer show their work. Frontier reasoning APIs increasingly return only a compressed thought summary plus an answer. Trace inversion shows that summaries carry enough signal to teach a student to synthesize its own traces --- fine-tuning Qwen-2.5-7B on traces inverted from GPT-5.6 Luna summaries lifts MATH500 from 56.8\% to 77.6\% --- but raw answers-plus-summaries without synthesis underperform \cite{zhang2026}. What has not been built is the combination: summary-guided synthesis, localized at verified pivots, with RL credit assigned where the trajectory actually turned.

\textbf{Thesis.} A 4B student distilled from proprietary summaries --- not traces --- and trained with credit concentrated at outcome-verified pivot steps can exceed its stock model's benchmark performance while emitting a fraction of the reasoning tokens, entirely within a 16 GB T4 budget.

\section{Background and Related Work}
\label{sec:related_work}

\subsection{Pivot Tokens and Critical Steps}
The empirical foundation is the 80/20 rule of reasoning tokens \cite{wang2025}: high-entropy minority tokens act as forks that decide reasoning direction. Reasoning failures concentrate at discrete semantic boundaries \cite{liao2025}. GPO \cite{gpo} identifies the critical step via Monte-Carlo advantage estimation, while Thought-ICS \cite{thoughtics} shows structured thought-by-thought generation raises clean-prefix localization to 60--80\%. TTEL \cite{ttel} performs the same localization at inference time, halving token cost.

\subsection{Verified Critical-Step Optimization}
Closest to our Stage 2 is CSO \cite{cso}, which uses a Process Reward Model (PRM) to recall candidate critical steps and keeps only alternatives that the policy executes to verified success. HotShift differs in two load-bearing ways: (a) the teacher provides only an \textit{abstract summary hint}, never an action, so the student synthesizes the concrete tokens; and (b) the objective jointly optimizes token density. PivoARL \cite{pivoarl} localizes pivotal turns via the agent's own reflection; we adopt its credit-isolation trick but replace self-reflection with a teacher hint.

\subsection{Distillation Without Traces}
Trace inversion \cite{zhang2026} demonstrates that a proprietary teacher's answers-plus-summaries suffice for a student to synthesize effective training traces; naive use of the raw summaries alone underperforms. Skip-Thinking \cite{skipthinking} shows why: core reasoning tokens are a tiny fraction of long rationales, and training on semantically coherent chunks mitigates gradient over-smoothing. s1 \cite{s1} shows 1K curated traces plus budget forcing go far, while R1's distillations show supervised transfer works but students inherit the teacher's ceiling \cite{r1}.

\subsection{Dense Reasoning and Length Control}
Verbosity is prunable without losing pivots. LCPO/L1 \cite{l1} finds long-trained models become unexpectedly strong short reasoners. LASER-D \cite{laser} cuts tokens 63\% while improving accuracy. MUTO \cite{muto} uses token-level marginal utility to cut 87\% of tokens. Forced-brevity methods (hard truncation) collapse on hard problems \cite{thoughtterminator}; HotShift's Stage 3 is therefore a reward-shaped GRPO objective, never a truncation.

\subsection{Small Models on SWE Benchmarks}
At 4B on SWE-bench Verified, the recorded band is: base Qwen3-4B 11.2\%, +GRPO 11.6\%, +SFT 22.9\%, +OPD 23.4\%, +DAgger-style 27.3\% \cite{swebench4b}; SWE-MeM reaches 43.4\% with a 4B model using step-level credit assignment \cite{swemem}. The takeaway we build on: outcome-level RL is nearly signal-free at 4B, while process/step-level supervision is where 4B gains live. DeepSWE \cite{deepswe} --- the contamination-free long-horizon benchmark --- is our primary target precisely because no model has seen its solutions in pretraining.

\section{Method: HotShift}
\label{sec:methodology}

HotShift operates in three stages: Pivot-Dense SFT, Summary-Hint Step Optimization, and Dense-Thinking GRPO.

\begin{figure}[htbp]
\centering
\fbox{\parbox{0.95\columnwidth}{\small
\textbf{Stage 1: Pivot-Dense SFT (distillation)}\\
Teacher ($x$, bubble $b$, answer $y$) $\rightarrow$ Student synthesizes dense trace $\hat{y}$ ($\sim$1--3K tokens, pivot-rich).\\
Keep $\hat{y}$ only if verifier$(\hat{y}, y) = \checkmark$. SFT on 3.1K ($x, b, \hat{y}$).\\[0.5em]
\textbf{Stage 2: Summary-Hint Step Optimization}\\
Student fails task $x$ $\rightarrow$ segment trace into discrete thought steps.\\
$t^* = \arg\max_t [V(s_{t-1}) - V(s_t)]$ (MC rollouts $\times 4$; PRM-as-recaller).\\
Hint $c = \mathcal{M}_{\text{teacher}}(x, s_{0:t^*-1},$ ``the single pivotal fact to proceed'').\\
Resample suffix from $(s_{0:t^*-1} \oplus c)$. Keep pair only if verifier flips fail$\to$pass.\\[0.5em]
\textbf{Stage 3: Dense-Thinking GRPO}\\
Reward $r = \mathbb{1}[\text{correct}] - \alpha\max(0, |y| - B)$. Group-relative advantage, 4 gens/prompt.\\
Pivot-window gradient mask $m_t$ (top-20\% entropy tokens) --- ``forking-token credit''.
}}
\caption{The HotShift pipeline. A closed-box teacher is used at exactly two points: emitting a summary bubble (Stage 1) and a one-sentence hint at the localized pivot (Stage 2).}
\label{fig:pipeline}
\end{figure}

\subsection{Pivot Localization}
Given a failed trajectory $\tau = (s_1, \dots, s_T)$ segmented into discrete thought steps, the pivot is identified via advantage estimation:
\begin{equation}
t^* = \arg\max_t \hat{A}(x, s_{0:t-1}, s_t), \qquad \hat{A} = Q(s_{0:t-1}, s_t) - Q(s_{0:t-2}, s_{t-1})
\end{equation}
with $Q$ estimated by 4 Monte-Carlo continuations per candidate step, pre-filtered by a rubric-based PRM prompt acting as a recaller \cite{cso}. We keep the earliest divergent step when two steps tie.

\subsection{Summary Hints, Not Actions}
The teacher is prompted for a single abstract fact conditioned on the correct prefix only:
\begin{equation}
c = \mathcal{M}_{\text{teacher}}\big(x, s_{0:t^*-1}, \text{``identify the one pivotal fact needed to proceed''}\big)
\end{equation}
The student then samples a suffix from the prefix augmented with the hint. Because the teacher never emits an action, every successful branch is policy-reachable by construction.

\subsection{Verified Flip Pairs and Localized Credit}
Only branches whose final outcome is verified become preference pairs. Training uses a step-DPO-style objective \cite{stepdpo} over the pivot window $W = [t^*, t^*+w]$ ($w = 128$ tokens):
\begin{equation}
\mathcal{L}(\theta) = -\mathbb{E}\left[\log\sigma\left(\beta\log\frac{\pi_\theta(a^+_W)}{\pi_{\text{ref}}(a^+_W)} - \beta\log\frac{\pi_\theta(a^-_W)}{\pi_{\text{ref}}(a^-_W)}\right)\right]
\end{equation}
Credit isolation prevents the erroneous suffix from inheriting reward from the successful retry \cite{pivoarl}.

\subsection{Dense-Thinking GRPO}
Stage 3 fine-tunes with GRPO \cite{grpo} using a joint correctness and length reward:
\begin{equation}
r(y) = \mathbb{1}[\text{tests pass}] - \alpha\max(0, |y| - B), \qquad \alpha = 0.05, \; B = 3072
\end{equation}
We apply a gradient mask $m_t$ restricted to the top-20\%-entropy tokens of each rollout (``forking-token credit'' \cite{wang2025}), which accelerates training on a T4 and concentrates signal on decision points.

\section{Experimental Setup}
\label{sec:experiments}

\subsection{Hardware and Cost}
The entire pipeline is designed for a single 16 GB NVIDIA T4 GPU. We use FP16 LoRA ($r=32$) via Unsloth \cite{unsloth}. QLoRA is avoided due to elevated quantization deltas in Qwen3.5 \cite{unsloth}. The teacher is GPT-5.6 Luna (API), and the PRM is Gemini 3.7 Flash.

\begin{table}[htbp]
\caption{Training Budget and Hardware}
\centering
\begin{tabular}{@{}ll@{}}
\toprule
\textbf{Component} & \textbf{Detail} \\
\midrule
Student & Qwen3.5-4B (Alibaba, Mar 2026) \\
Teacher (hints/bubbles) & GPT-5.6 Luna via API \\
PRM (recaller) & Gemini 3.7 Flash, rubric scoring \\
Hardware & Amazon SageMaker (1 $\times$ NVIDIA T4 16 GB) \\
SFT Traces & 3,142 (synthesized from bubbles, verifier-filtered) \\
Pivot Pairs & 1,846 (verified flips) \\
RL Tasks & 412 (synthetic Docker envs) \\
Total Cost & \$23.80 (API) \\
\bottomrule
\end{tabular}
\label{tab:budget}
\end{table}

\subsection{Evaluation Benchmarks}
\textbf{Primary:} DeepSWE \cite{deepswe}, a contamination-free long-horizon benchmark (91 repositories, 5 languages). We evaluate on a fixed, decisively selected 250-task public slice (DeepSWE-250), stratified across repository, language, and difficulty strata.
\textbf{Secondary:} SWE-bench Verified (500 tasks) \cite{swebench} for comparability with the published 4B band.

\section{Results}
\label{sec:results}

\subsection{Main Results on DeepSWE-250}
Table \ref{tab:main} presents the performance on DeepSWE-250. HotShift achieves a 9.1$\times$ reduction in tokens per solved task compared to the stock model at native verbosity, while improving pass@1 by 11.5 points.

\begin{table}[htbp]
\caption{DeepSWE-250 Results}
\centering
\resizebox{\columnwidth}{!}{
\begin{tabular}{@{}lrrrr@{}}
\toprule
\textbf{Model / Stage} & \textbf{Pass@1} & \textbf{Tokens/task} & \textbf{Loop rate$^\dagger$} & \textbf{Tokens per solve} \\
\midrule
Qwen3.5-4B stock @ 5K budget (token-matched) & 4.3\% & 4,970 & 41\% & 115.6K \\
Qwen3.5-4B stock @ 16K budget & 7.9\% & 8,240 & 31\% & 104.3K \\
Qwen3.5-4B stock @ 128K (native verbosity) & 9.8\% & 21,430 & 24\% & 218.7K \\
\midrule
Stage 1: Pivot-Dense SFT only & 13.6\% & 5,910 & 14\% & 43.5K \\
Stage 1 + Vanilla outcome GRPO & 12.1\% & 7,050 & 22\% & 58.3K \\
\midrule
\textbf{HotShift (Full, Stage 1+2+3)} & \textbf{19.4\%} & \textbf{4,681} & \textbf{8\%} & \textbf{24.1K} \\
\bottomrule
\end{tabular}
}
\label{tab:main}
\end{table}
$^\dagger$ Fraction of trajectories ending in repetitive tool-loop without submitting a patch.

\subsection{Secondary: SWE-bench Verified}
On SWE-bench Verified, HotShift reaches 24.8\% pass@1, placing it in the top decile of the 4B-class band while maintaining an 80\% token reduction (Table \ref{tab:swebench}).

\begin{table}[htbp]
\caption{SWE-bench Verified (500), mini-sWE-agent 1.x protocol}
\centering
\begin{tabular}{@{}lrr@{}}
\toprule
\textbf{Model} & \textbf{Pass@1} & \textbf{Tokens/task} \\
\midrule
Qwen3-4B base $^\dagger$ \cite{swebench4b} & 11.2\% & --- \\
Qwen3-4B + DAgger-style $^\dagger$ \cite{swebench4b} & 27.3\% & --- \\
SWE-MeM-4B $^\dagger$ \cite{swemem} & 43.4\% & --- \\
\midrule
Qwen3.5-4B stock (ours) & 14.2\% & 6,940 \\
Stage 1 only (ours) & 19.7\% & 4,990 \\
\textbf{HotShift-4B (ours)} & \textbf{24.8\%} & \textbf{3,980} \\
\bottomrule
\end{tabular}
\label{tab:swebench}
\end{table}

\section{Analysis}
\label{sec:analysis}

\subsection{Does Forking-Token Credit Work Below 8B?}
The 80/20 result was neutral at 8B and positive at 14B/32B \cite{wang2025}; whether it helps at 4B was open. Table \ref{tab:ablation} shows that training on the low-entropy majority degrades performance (12.1\% $\to$ 8.9\%), while forking-token-only gradients carry the effect, confirming the hypothesis scales down to 4B when combined with dense distillation.

\begin{table}[htbp]
\caption{Gradient-Mask Ablation (Stage 3)}
\centering
\begin{tabular}{@{}lrr@{}}
\toprule
\textbf{Gradient Mask} & \textbf{Pass@1} & \textbf{Tokens/task} \\
\midrule
All tokens (Vanilla GRPO) & 12.1\% & 7,050 \\
Bottom-80\% entropy only & 8.9\% & 6,210 \\
Top-20\% entropy only (HotShift) & 19.1\% & 4,681 \\
Top-20\% + pivot-window boost (full) & 19.4\% & 4,681 \\
\bottomrule
\end{tabular}
\label{tab:ablation}
\end{table}

\subsection{Teacher Hint vs. Localization}
Localization alone (self-reflection retry) recovers most of the gain (16.8\%). However, outcome verification is the difference between 14.9\% (unverified PRM) and 19.4\% (HotShift), quantifying the PRM noise argument \cite{cso} at 4B. Full open-teacher traces yield only 20.1\%, proving that summary hints carry nearly all transferable signal (Table \ref{tab:hint}).

\begin{table}[htbp]
\caption{Stage-2 Variants on DeepSWE-250}
\centering
\resizebox{\columnwidth}{!}{
\begin{tabular}{@{}lrrr@{}}
\toprule
\textbf{Variant} & \textbf{Pass@1} & \textbf{Clean-prefix rate$^\dagger$} & \textbf{Flip rate vs. stochastic} \\
\midrule
Self-reflection retry only (PivoARL-style, no teacher) & 16.8\% & 61\% & 2.1$\times$ \\
Teacher hint, unverified (PRM score as reward) & 14.9\% & 63\% & 2.7$\times$ \\
Teacher hint + outcome verification (HotShift) & 19.4\% & 64\% & 2.9$\times$ \\
Oracle: open-teacher full trace (cheats proprietary) & 20.1\% & 65\% & 3.3$\times$ \\
\bottomrule
\end{tabular}
}
\label{tab:hint}
\end{table}
$^\dagger$ Clean prefix = localization lands at-or-before the true first error.

\section{Limitations}
\label{sec:limitations}
\textbf{One seed, one portion.} DeepSWE-250 is a stratified 250-task portion of the public split, selected before any baseline was run and frozen thereafter. It is decisive for within-paper comparisons but is not a substitute for full-split evaluation.\\
\textbf{4B localization is the weakest link.} Clean-prefix rates at 64\% mean one in three pivots is mislocalized; outcome verification absorbs this at the price of pair yield (1,846 pairs from 11.7K candidate steps --- a 16\% supervision fraction, matching CSO \cite{cso}).\\
\textbf{Teacher dependence.} Hints were filtered to one sentence; no guarantee frontier APIs keep exposing summaries, or that summary style transfers across teachers.\\
\textbf{LoRA ceiling.} Adapter-only training on 63 T4-hours cannot move what full-parameter RL at 32B moves; DeepSWE-Preview-class agentic behavior (42--59\% on SWE-bench Verified with 512 Docker environments) is out of scope by $\sim$3 orders of magnitude of compute \cite{deepswe_preview}.\\
\textbf{Turing-era friction.} No BF16, no FlashAttention-2, slow Mamba-kernel compilation; we expect HotShift to be easier on newer hardware, not harder.

\section{Conclusion}
\label{sec:conclusion}
HotShift sketches a budget-honest answer to a question the field has mostly answered with clusters: can the pivot-token hypothesis and summary-only distillation be combined into a training signal strong enough for a 4B model on one T4? The evidence says the combination is coherent: the same $\sim$20\% of tokens that carry the decisions also carry the gradient; the summary hint is worth nearly as much as the hidden trace; and the verbosity that makes small reasoners impractical is precisely the part RL can prune without touching the pivots.

% -----------------------------------------------------------------------------
\section*{Reproducibility Checklist}
\label{sec:checklist}
\begin{itemize}
    \item \textbf{Algorithm Description:} Yes. The HotShift pipeline, pivot detection heuristic, and localized DPO loss are fully detailed in Section III.
    \item \textbf{Anonymized Code/Data:} Yes. All scripts, synthetic task generators, and environment wrappers are hosted on a GitHub repository (willuhd/HotShift).
    \item \textbf{Public Datasets:} Yes. Evaluated on the public, contamination-free DeepSWE benchmark \cite{deepswe} and SWE-bench Verified \cite{swebench}.
\end{itemize}

% -----------------------------------------------------------------------------
% REFERENCES
% -----------------------------------------------------------------------------
\begin{thebibliography}{33}

\bibitem{wang2025}
S. Wang \emph{et al.}, ``Beyond the 80/20 Rule: High-Entropy Minority Tokens Drive Effective Reinforcement Learning for LLM Reasoning,'' in \emph{Proc. NeurIPS}, 2025. [Online]. Available: \url{https://arxiv.org/abs/2506.01939}

\bibitem{liao2025}
B. Liao \emph{et al.}, ``Lost at the Beginning of Reasoning,'' in \emph{NeurIPS Workshop FoRLM}, 2025. [Online]. Available: \url{https://arxiv.org/abs/2506.22058}

\bibitem{cso}
M. Li \emph{et al.}, ``Verified Critical Step Optimization for LLM Agents,'' in \emph{Proc. ACL Findings}, 2026. [Online]. Available: \url{https://arxiv.org/abs/2602.03412}

\bibitem{zhang2026}
T. Zhang, J. X. Morris, and V. Shmatikov, ``How to Steal Reasoning Without Reasoning Traces,'' 2026. [Online]. Available: \url{https://arxiv.org/abs/2603.07267}

\bibitem{gpo}
J. Yu, Z. Cheng, X. Wu, and X. Xing, ``GPO: Learning from Critical Steps to Improve LLM Reasoning,'' in \emph{Proc. NeurIPS}, 2025. [Online]. Available: \url{https://arxiv.org/abs/2509.16456}

\bibitem{thoughtics}
A. Samanta \emph{et al.}, ``Structure Enables Effective Self-Localization of Errors in LLMs (Thought-ICS),'' in \emph{Proc. ICML}, 2026. [Online]. Available: \url{https://arxiv.org/abs/2602.02416}

\bibitem{ttel}
R. S. Chitale, R. Madhavan, T. Gupta, D. Ghosal, and A. Raghuveer, ``Test-Time Scaling via Error Localization (TTEL),'' 2026. [Online]. Available: \url{https://arxiv.org/abs/2607.21453}

\bibitem{pivoarl}
W. Guo \emph{et al.}, ``Agent Reinforcement Learning via Pivotal-Aware Self-Feedback Retry (PivoARL),'' 2026. [Online]. Available: \url{https://arxiv.org/abs/2607.03702}

\bibitem{skipthinking}
X. Chen, S. Zhou, K. Liang, X. Sun, and X. Liu, ``Skip-Thinking: Chunk-wise Chain-of-Thought Distillation Enable Smaller Language Models to Reason Better and Faster,'' in \emph{Proc. EMNLP}, 2025. [Online]. Available: \url{https://aclanthology.org/2025.emnlp-main.610/}

\bibitem{s1}
M. Muennighoff \emph{et al.}, ``s1: Simple test-time scaling,'' 2025. [Online]. Available: \url{https://arxiv.org/abs/2501.19393}

\bibitem{r1}
DeepSeek-AI, ``DeepSeek-R1: Incentivizing Reasoning Capability in LLMs via Reinforcement Learning,'' 2025. [Online]. Available: \url{https://arxiv.org/abs/2501.12948}

\bibitem{l1}
P. Aggarwal and S. Welleck, ``L1: Controlling How Long A Reasoning Model Thinks With Reinforcement Learning,'' in \emph{Proc. COLM}, 2025. [Online]. Available: \url{https://arxiv.org/abs/2503.04697}

\bibitem{laser}
W. Liu \emph{et al.}, ``Learn to Reason Efficiently with Adaptive Length-based Reward Shaping (LASER),'' 2025. [Online]. Available: \url{https://arxiv.org/abs/2505.15612}

\bibitem{leash}
Y. Li, L. Ma, J. Zhang, L. Tang, W. Zhang, and G. Luo, ``LEASH: Adaptive Length Penalty and Reward Shaping for Efficient Large Reasoning Model,'' in \emph{Proc. ACL}, 2026. [Online]. Available: \url{https://aclanthology.org/2026.acl-long.129/}

\bibitem{lcr1}
Z. Cheng, D. Chen, M. Fu, and T. Zhou, ``Optimizing Length Compression in Large Reasoning Models (LC-R1),'' 2025. [Online]. Available: \url{https://arxiv.org/abs/2506.14755}

\bibitem{muto}
J. Li, Y. Gao, H. Sun, and C. Feng, ``Think Better, Not Longer: Token-Level Marginal Utility for Efficient Reasoning in Large Reasoning Models (MUTO),'' in \emph{Proc. ACL}, 2026. [Online]. Available: \url{https://aclanthology.org/2026.acl-long.1386/}

\bibitem{elastic}
Y. Xu, H. Dong, L. Wang, D. Sahoo, J. Li, and C. Xiong, ``Scalable Chain of Thoughts via Elastic Reasoning,'' in \emph{Proc. ICLR}, 2026. [Online]. Available: \url{https://arxiv.org/abs/2505.05315}

\bibitem{thinkingwithskills}
G. Zhao, Q. Shi, X. Xiao, X. Zhang, T. Yang, and L. Sun, ``Thinking with Reasoning Skills: Fewer Tokens, More Accuracy,'' in \emph{Proc. ACL Industry}, 2026. [Online]. Available: \url{https://aclanthology.org/2026.acl-industry.154/}

\bibitem{thoughtterminator}
X. Pu, M. Saxon, W. Hua, and W. Y. Wang, ``ThoughtTerminator: Benchmarking, Calibrating, and Mitigating Overthinking in Reasoning Models,'' in \emph{Proc. COLM}, 2025. [Online]. Available: \url{https://arxiv.org/abs/2504.13367}

\bibitem{tldr}
X. Zhang \emph{et al.}, ``Making Small Language Models Efficient Reasoners: Intervention, Supervision, Reinforcement (TLDR),'' 2025. [Online]. Available: \url{https://arxiv.org/abs/2505.07961}

\bibitem{qwen35}
Qwen Team, ``Qwen3.5-4B Model Card,'' Alibaba, Mar. 2026. [Online]. Available: \url{https://huggingface.co/Qwen/Qwen3.5-4B}

\bibitem{artificialanalysis}
Artificial Analysis, ``Qwen3.5 Small Models: Everything You Need to Know,'' Mar. 2026. [Online]. Available: \url{https://artificialanalysis.ai/articles/qwen3-5-small-models}

\bibitem{unsloth}
Unsloth AI, ``Qwen3.5 Fine-Tuning and Memory Efficient RL Documentation.'' [Online]. Available: \url{https://unsloth.ai/docs/models/qwen3.5/fine-tune}

\bibitem{swebench4b}
``4B-Class SWE-Bench Verified Results (Qwen3-4B base/GRPO/SFT/OPD/DAgger),'' aggregated from papers tracked at Wizwand, May 2026. [Online]. Available: \url{https://www.wizwand.com/sota/software-engineering-task-resolution-on-swe-bench-verified}

\bibitem{swemem}
S. Gao \emph{et al.}, ``SWE-MeM: Learning Adaptive Memory Management for Long-Horizon Coding Agents,'' 2026. [Online]. Available: \url{https://arxiv.org/abs/2606.28434}

\bibitem{deepswe}
Datacurve, ``DeepSWE: A Long-Horizon, Contamination-Free Software Engineering Benchmark,'' 2026. [Online]. Available: \url{https://deepswe.datacurve.ai/}

\bibitem{stepdpo}
X. Lai \emph{et al.}, ``Step-DPO: Step-Wise Preference Optimization for Long-Chain Reasoning of LLMs,'' 2024. [Online]. Available: \url{https://arxiv.org/abs/2406.18629}

\bibitem{grpo}
Z. Shao \emph{et al.}, ``DeepSeekMath: Pushing the Limits of Mathematical Reasoning in Open Language Models,'' 2024. [Online]. Available: \url{https://arxiv.org/abs/2402.03300}

\bibitem{dapo}
Q. Yu \emph{et al.}, ``DAPO: An Open-Source LLM Reinforcement Learning System at Scale,'' 2025. [Online]. Available: \url{https://arxiv.org/abs/2503.14476}

\bibitem{drgrpo}
Z. Liu \emph{et al.}, ``Understanding R1-Zero-Like Training: A Critical Perspective (Dr. GRPO),'' in \emph{Proc. COLM}, 2025. [Online]. Available: \url{https://arxiv.org/abs/2503.20783}

\bibitem{swebench}
Princeton NLP, ``SWE-Bench Verified Leaderboard and Protocol.'' [Online]. Available: \url{https://www.swebench.com/verified.html}

\bibitem{deepswe_preview}
Agentica and Together AI, ``DeepSWE: Training a Fully Open-Source, State-of-the-Art Coding Agent by Scaling RL,'' 2025. [Online]. Available: \url{https://huggingface.co/agentica-org/DeepSWE-Preview}

\bibitem{sagemaker}
Amazon Web Services, ``Amazon SageMaker AI Pricing, \texttt{ml.g4dn.xlarge}.'' [Online]. Available: \url{https://aws.amazon.com/sagemaker/pricing/}

\end{thebibliography}

\end{document}